\scanpage{141}
\begin{center}- 87 -\end{center}
\margin{n° 239}
Supposons par exemple que $\mathcal G$ soit fin, et, avec les notations précédentes, soit $\theta$ un endomorphisme de $\mathcal G$ vérifiant $(\mathrm{FM}''_{\mathrm{II}})$ ; $1\otimes\theta$ est alors un endomorphisme de $\mathcal F\otimes\mathcal G$ vérifiant la même condition, d'où la proposition.

PROPOSITION 8. - \uline{Tout faisceau de} $\mathcal A$-\uline{modules} $\mathcal F$ \uline{est isomorphe à un sous-faisceau d'un faisceau fin de} $\mathcal A$-\uline{modules}.

Désignant encore par $\mathbb Z$ le faisceau constant d'anneaux $X\times\mathbb Z$, considérons $\mathcal F$ comme un faisceau de $\mathbb Z$-modules à droite (chaque fibre étant donc un $(\mathcal A(x),\mathbb Z)$-bimodule). Soit $\unsure{\mathcal I|\mathcal J}$ le faisceau (d'anneaux) des germes d'applications de $X$ dans $\mathbb Z$ ; on a vu plus haut (n°4,\uline{Exemple} 1) que $\unsure{\mathcal I|\mathcal J}$ est un faisceau mou d'anneaux commutatifs, donc aussi (prop.5) un faisceau fin quand on le considère comme faisceau de $\unsure{\mathcal I|\mathcal J}$-modules ; a fortiori, $\unsure{\mathcal I|\mathcal J}$ est un faisceau fin de $\mathbb Z$-modules. D'après la prop.7, le faisceau $\mathcal F\otimes_{\mathbb Z}\unsure{\mathcal I|\mathcal J}$ est donc un faisceau fin de $\mathbb Z$-modules, et il est muni d'une structure de faisceau de $\mathcal A$-modules à gauche. Tout revient à prouver que l'homomorphisme $\mathcal A$-linéaire $f\longrightarrow f\otimes1$ de $\mathcal F$ dans $\mathcal F\otimes\unsure{\mathcal I|\mathcal J}$ est une \uline{injection} (§ 1, prop.6) ; cela sera le cas si la restriction de cette application à $\mathcal F(x)$ (qui est une application $\mathcal A(x)$-linéaire de $\mathcal F(x)$ dans $\mathcal F(x)\otimes\unsure{\mathcal I|\mathcal J}(x)$) est injective. Or, il est immédiat que le groupe abélien $\unsure{\mathcal I|\mathcal J}(x)$ est \uline{sans torsion} ; le résultat est donc conséquence du lemme algébrique suivant :

LEMME 2. - \uline{Soient} $G,H$ \uline{deux groupes abéliens}, $H$ \uline{étant sans torsion}. \uline{Si} $a\in H$ \uline{est} $\ne0$ \uline{et tel que la relation} $mb=a$ ($m$ \uline{entier}) \uline{entraîne} $b=\pm a$, \uline{l'application} $x\longrightarrow x\otimes a$ \uline{de} $G$ \uline{dans} $G\otimes H$ \uline{est injective}.

Il suffit de le démontrer lorsque $H$ est \uline{de type fini} (\uline{Alg.}, chap.III, § 1, prop.8) ; mais alors, comme $H$ est sans torsion, on sait
\mspage{140}

